% GENERATED FILE. Edit canonical lessons, then run npm run generate:books.
% canonical-source-hash: fnv1a64:8189dbd0f7189ffa
% canonical-lessons: MW-C04-hear-paani, MW-W04-pa, MW-W04-nna, MW-W04-ii-matra, MW-C04-paani, MW-C04-practice

\chapter{Paani: Hear Water, Then Write It Sign by Sign}
\label{ch:mw-paani}
\hlchaptermodality{4}

\begin{chapteropening}
\textbf{By the end of this chapter:} \emph{I can recognise, hear, say, read, and independently write \mw{पाणी} as the Marwadi word for water.}

A source-attested everyday noun secured by ear before three new signs are isolated, assembled, and retrieved independently in all four skills.
\end{chapteropening}

\section[pāṇī]{Hear \emph{pāṇī} — meaning before spelling}

\label{lesson:MW-C04-hear-paani}



\begin{quote}
Give the Chapter 3 four-skill response once: hear \emph{hā(n) sā} and identify a respectful yes; say it; read \textbf{\mw{हां} \mw{सा}}; then cover and write it. Say \emph{ābhār} for formal thanks. Now set writing aside: today's new word begins with the ear and voice.
\end{quote}

\subsection*{You'll want to know — one useful word}
\begin{quote}
\emph{pāṇī} — water
\end{quote}

Listen for two long vowels: \emph{pā-ṇī}. The middle sound is made with the tongue curled back a little. Hear the whole word twice, then say it once. Do not try to spell it yet; the next three tiny lessons will give the hand one piece at a time.

\subsection*{Your turn — meaning without letters}
Imagine two cups: one holds water, one holds tea. Hear \emph{pāṇī} and point to the water. Change the order and do it again. Then say \emph{pāṇī} while imagining water.

\begin{tcolorbox}[breakable,colback=teal!4,colframe=teal!35!black,title={Before you move on}]
Close the page. Say \emph{pāṇī} once, then give its meaning.

Source: \href{https://nepal.sil.org/resources/archives/63768}{SIL International, \emph{Marwari–English Dictionary} (2015)}.
\end{tcolorbox}
\section[pa]{\mw{प} — \emph{pa}, the first sign}

\label{lesson:MW-W04-pa}



\begin{quote}
Write \textbf{\mw{हां}} from memory, checking the first sign and the dot. With no new spelling in view, say the Marwadi word for water. Then write the familiar long-\emph{ā} mark \textbf{\mw{ा}} once.
\end{quote}

\subsection*{Script — one new consonant}
\begin{quote}
\textbf{\mw{प}} — \emph{pa}
\end{quote}

Close the lips, release them without a strong puff, and say \emph{pa}. Today the hand meets only \textbf{\mw{प}}; the complete word stays hidden.

\subsection*{Writing — observe, trace, copy}
1. look at \textbf{\mw{प}} for three seconds 2. finger-trace its left stem, bowl, and headline 3. copy \textbf{\mw{प}} twice 4. say \emph{pa} after each copy

\begin{tcolorbox}[breakable,colback=teal!4,colframe=teal!35!black,title={Before you move on}]
Hide the model for five seconds. Write \textbf{\mw{प}}, then add \textbf{\mw{ा}} to make \textbf{\mw{पा}}.
\end{tcolorbox}
\section[ṇa]{\mw{ण} — curled-back \emph{ṇa}}

\label{lesson:MW-W04-nna}



\begin{quote}
Write \textbf{\mw{हां}} as one complete known word. Say \emph{pāṇī} without looking at its spelling, then write \textbf{\mw{पा}} from memory.
\end{quote}

\subsection*{Script — one new consonant}
\begin{quote}
\textbf{\mw{ण}} — \emph{ṇa}
\end{quote}

Curl the tongue tip slightly back for \emph{ṇ}. Do not replace it with the dental \emph{n} heard in the Hindi cognate. The Marwadi dictionary spelling in this chapter uses \textbf{\mw{ण}}, so the eye, voice, and hand keep that contrast clear from the beginning.

\subsection*{Writing — observe, trace, copy}
1. look at \textbf{\mw{ण}} for three seconds 2. finger-trace the lower loop, upright, and headline 3. copy \textbf{\mw{ण}} twice 4. say \emph{ṇa} after each copy

\begin{tcolorbox}[breakable,colback=teal!4,colframe=teal!35!black,title={Before you move on}]
Cover both models. Write \textbf{\mw{प}}, then \textbf{\mw{ण}}, with a clear space between them.
\end{tcolorbox}
\section[long ī mark]{\mw{ी} — the long-\emph{ī} mark}

\label{lesson:MW-W04-ii-matra}



\begin{quote}
Hear “Will you come?” and answer yes respectfully. Then say \emph{pāṇī}. Write \textbf{\mw{प}}, then \textbf{\mw{ण}}, from memory and name the curled-back sound.
\end{quote}

\subsection*{Script — one new vowel mark}
\begin{quote}
\textbf{\mw{ी}} — long \emph{ī}
\end{quote}

Like familiar \textbf{\mw{ा}}, this mark cannot stand alone in a word. It attaches to a consonant. Add it to \textbf{\mw{ण}} and read \textbf{\mw{णी}}, \emph{ṇī}.

\subsection*{Writing — attach only the mark}
1. copy \textbf{\mw{ण}} once 2. look at \textbf{\mw{णी}} 3. add \textbf{\mw{ी}} on the right of your copy 4. compare the joined form, then say \emph{ṇī}

\begin{tcolorbox}[breakable,colback=teal!4,colframe=teal!35!black,title={Before you move on}]
Hide \textbf{\mw{णी}} for five seconds. Write it once, then uncover and compare.
\end{tcolorbox}
\section[pāṇī]{\mw{पाणी} — join the word you already know by ear}

\label{lesson:MW-C04-paani}



\begin{quote}
Write the first sign of \textbf{\mw{राम}} from memory. Then say the Marwadi word for water and write \textbf{\mw{पा}}, then \textbf{\mw{णी}}, as two separate pieces. Check each piece before going on.
\end{quote}

\subsection*{The exchange — two familiar pieces}
\begin{quote}
\textbf{\mw{पा}} + \textbf{\mw{णी}} $\to$ \textbf{\mw{पाणी}} — \emph{pāṇī} — water
\end{quote}

Join the pieces without adding or changing a sign. Keep \textbf{\mw{ण}}, the curled-back sound that distinguishes this source-attested Marwadi spelling from its Hindi cognate.

\begin{cousinweb}
The word belongs to an old Indo-Aryan family. Sanskrit \emph{pānīya} meant something for drinking; related regional forms travelled through Middle Indo-Aryan into modern Marwadi \textbf{\mw{पाणी}}. The history is a memory hook, not a new spelling task.
\end{cousinweb}

\subsection*{Writing — guided, then delayed}
1. copy \textbf{\mw{पाणी}} once while saying \emph{pā-ṇī} 2. point to \textbf{\mw{पा}}, then \textbf{\mw{णी}} 3. hide the model for ten seconds 4. write the complete word from memory 5. uncover and compare \textbf{\mw{प}}, \textbf{\mw{ा}}, \textbf{\mw{ण}}, and \textbf{\mw{ी}} separately

\begin{tcolorbox}[breakable,colback=teal!4,colframe=teal!35!black,title={Before you move on}]
Say \textbf{\mw{पाणी}} once without looking, then write it once with the model covered.

Sources: \href{https://nepal.sil.org/resources/archives/63768}{SIL International, \emph{Marwari–English Dictionary} (2015)} and \href{https://dsal.uchicago.edu/dictionaries/soas/}{Turner, \emph{A Comparative Dictionary of the Indo-Aryan Languages}, entry 8082}.
\end{tcolorbox}
\section[Practice]{Practice — hear, say, read, and write water}

\label{lesson:MW-C04-practice}



\begin{quote}
Give the older respectful yes in four quick steps: hear it and name its meaning; say it; read \textbf{\mw{हां} \mw{सा}}; then cover and write it. Now say the three meanings in English as you hear \emph{ābhār}, \emph{hā(n) sā}, and \emph{pāṇī}. Write \textbf{\mw{पा}}, then \textbf{\mw{ण}} and \textbf{\mw{ी}} separately from memory.
\end{quote}

\subsection*{Your turn — four independent checks}
1. \textbf{Listen:} hear \emph{pāṇī} and identify water, without text. 2. \textbf{Speak:} see a picture or imagine a glass of water and say \emph{pāṇī}, without notes. 3. \textbf{Read:} choose \textbf{\mw{पाणी}} from \textbf{\mw{आभार} · \mw{हां} · \mw{पाणी}} and give its meaning. 4. \textbf{Write:} hear \emph{pāṇī}, wait fifteen seconds, then write \textbf{\mw{पाणी}} with no model.

Score each line separately. Repeat only the missed line; strong speaking does not compensate for a missing \textbf{\mw{ण}} or \textbf{\mw{ी}} in writing.

\begin{grammarlens}[title={What you've built — and what remains}]
You can retrieve one useful word in all four skills. You cannot yet request water in a complete Marwadi sentence; that needs pronouns, a request form, and register choices taught in later small steps.
\end{grammarlens}

\begin{tcolorbox}[breakable,colback=teal!4,colframe=teal!35!black,title={Before you move on}]
If one sign stalled, return only to its two-minute sign lesson before repeating the complete dictation.
\end{tcolorbox}
